% !TeX root = elegantnote-en.tex
\section{Reactions of Aromatic Compounds}

\subsection{Electrophilic Aromatic Substitution 20260326}

\begin{enumerate}
    \item Although benzene’s pi electrons are in a stable aromatic system, they are available to attack a strong electrophile to give a carbocation.
    \item This resonance-stabilized carbocation is called a \textbf{sigma complex} because the electrophile is joined to the benzene ring by a new sigma bond.
    \item The sigma complex (also called an arenium ion) is not aromatic because the \textit{sp}$^3$ hybrid carbon atom interrupts the ring of \textit{p} orbitals.
    \begin{figure}[h]
        \centering
        \includegraphics[width=1.0\textwidth]{17.1_Sigma_Complex.png}
        \caption{Sigma Complex}
        \label{17.1_Sigma_Complex}
    \end{figure}
\end{enumerate}

\begin{mechanism}
    Electrophilic Aromatic Substitution
    \begin{itemize}
        \item Step 1: Attack on the electrophile forms the sigma complex.
        \item Step 2: Loss of a proton regains aromaticity and gives the substitution product.
    \end{itemize}

    \begin{example} % no sequence
        Iodination of toluene
    \end{example}
\end{mechanism}

\subsection{Halogenation of Benzene 20260326}

\subsubsection{Bromination of Benzene}

\begin{enumerate}
    \item Bromination follows the general mechanism for electrophilic aromatic substitution.
    \item Bromine itself is not sufficiently electrophilic, and the formation of \ch{Br+} is difficult.
    \item A strong Lewis acid such as catalyzes the reaction by forming a complex with \ch{Br2} that reacts like \ch{Br+}.
\end{enumerate}

\begin{mechanism}
    Bromination of Benzene
\end{mechanism}

\subsubsection{Chlorination of Benzene}

\begin{enumerate}
    \item Chlorination of benzene works much like bromination, except that aluminum chloride (\ch{AlCl3}) is most often used as the Lewis acid catalyst.
\end{enumerate}

\subsubsection{Iodination of Benzene}

\begin{enumerate}
    \item Iodination of benzene requires an acidic oxidizing agent, such as nitric acid, to produce iodide cation.
\end{enumerate}

\begin{example}
    flourination
\end{example}

\subsection{Nitration of Benzene 20260326}

\begin{enumerate}
    \item Benzene reacts with hot, concentrated nitric acid to give nitrobenzene. A safer and more convenient procedure uses a mixture of nitric acid and sulfuric acid.
    \item Sulfuric acid is a catalyst, allowing nitration to take place more rapidly and at lower temperatures.
\end{enumerate}

\begin{mechanism}
    Nitration of Benzene
    \begin{itemize}
        \item Preliminary steps: Formation of the nitronium ion, \ch{NO2+}.
        \item Step 1: Attack on the electrophile forms the sigma complex.
        \item Step 2: Loss of a proton gives nitrobenzene.
    \end{itemize}
\end{mechanism}

\begin{enumerate}[resume]
    \item Aromatic nitro groups are easily reduced to amino (\chemfig{-NH_2}) groups by treatment with an active metal such as tin, zinc, or iron in dilute acid.
    \item Nitration followed by reduction is often the best method for adding an amino group to an aromatic ring.
\end{enumerate}

\subsection{Sulfonation of Benzene 20260326}

\subsubsection{Sulfonation}

\begin{enumerate}
    \item Sulfur trioxide (\ch{SO3}) is the electrophile.
    \item “Fuming sulfuric acid” is the common name for a solution of 7\% \ch{SO3} in \ch{H2SO4}.
\end{enumerate}

\begin{mechanism}
    Sulfonation of Benzene
\end{mechanism}

\subsubsection{Desulfonation}

\begin{enumerate}
    \item A sulfonic acid group may be removed from an aromatic ring by heating in dilute sulfuric acid.
    \item Steam is often used as a source of both water and heat for desulfonation.
\end{enumerate}

\subsubsection{Protonation of the Aromatic Ring; Hydrogen–Deuterium Exchange}

\subsection{Nitration of Toluene: The Effect of Alkyl Substitution 20260326}

\begin{enumerate}
    \item Toluene (methylbenzene) reacts with a mixture of nitric and sulfuric acids much like benzene does, but with some interesting differences:
    \begin{enumerate}
        \item Toluene reacts about 25 times faster than benzene under the same conditions. Toluene is activated toward electrophilic aromatic substitution, and the methyl group is an \textbf{activating group}.
        \item Nitration of toluene gives a mixture of products, primarily those resulting from substitution at the ortho and para positions. The methyl group of toluene is an \textbf{ortho, para-director}.
    \end{enumerate}
    \item In ortho or para substitution of toluene, the positive charge is spread over two secondary carbons and one tertiary ($ 3^{\circ} $) carbon (bearing the \ch{CH3} group).
    \item The sigma complex for meta substitution has its positive charge spread over three $ 2^{\circ} $ carbons. 
    \item Inductive effect: % position? 17-6?
    \begin{itemize}
        \item Electron-withdrawing: \chemfig{-F}, \chemfig{-Cl}, \chemfig{-Br}, \chemfig{-NO_2}, \chemfig{-O}, \chemfig{-N}
        \item Electron-donating: \chemfig{-CH_3}, \chemfig{-R}
    \end{itemize}
    Resonance effect:
    \begin{itemize}
        \item Electron-withdrawing: \chemfig{-NO_2}, \chemfig{-C~N}, \chemfig{-C=O}
        \item Electron-donating: \chemfig{-OR}, \chemfig{-NR_2}, \chemfig{-Cl}, \chemfig{-Br}
    \end{itemize}
    
\end{enumerate}

\subsection{Activating, Ortho, Para-Directing Substituents 20260331}

\subsubsection{Alkyl Groups}

\begin{enumerate}
    \item An alkyl group is an activating substituent, and it is ortho, para-directing.
    \item Substitution ortho or para to the alkyl group gives a transition state and an intermediate with the positive charge shared by the tertiary carbon atom.
    \item This effect is called inductive stabilization because the alkyl group donates electron density through the sigma bond joining it with the benzene ring.
    \item
\end{enumerate}

\begin{example} % 17-6
    Styrene (vinylbenzene) undergoes electrophilic aromatic substitution much faster than benzene, and the products are found to be primarily ortho- and para-substituted styrenes. A phenyl substituent should be ortho, para-directing like vinyl.
\end{example}

\subsubsection{Substituents with Nonbonding Electrons}

\begin{enumerate}
    \item Alkoxyl Groups
    \begin{enumerate}
        \item Anisole (methoxybenzene) undergoes nitration about 10,000 times faster than benzene.
        \item The methoxy group of anisole preferentially activates the ortho and para positions.
        \item The nonbonding electrons of an oxygen atom adjacent to a carbocation stabilize the positive charge through resonance.
        \item This type of stabilization is called \textbf{resonance stabilization}, and the oxygen atom is called \textbf{resonance-donating} or \textbf{pi-donating}.
        \item The resonance effect is much stronger than the inductive effect.
        \item A methoxy group is so strongly activating that anisole quickly brominates in water without a catalyst. In the presence of excess bromine, this reaction proceeds to the tribromide. This reaction can be controlled in low temerature and with a limited amount of bromine to give the single or double substituted product.
    \end{enumerate}
    \item Amine Groups
    \begin{enumerate}
        \item A nitrogen atom with a nonbonding pair of electrons serves as a powerful ortho, para director and activating group. 
        \item Aniline undergoes a fast bromination (without a catalyst) in bromine water to give the tribromide.
        \item Sodium bicarbonate is added to neutralize the HBr formed and to prevent protonation of the basic amino (\chemfig{-NH_2}) group, which would be a electron-withdrawing, meta-directing substituent.
        \item Nitrogen’s nonbonding electrons provide resonance stabilization to the sigma complex if attack takes place ortho or para to the position of the nitrogen atom.
    \end{enumerate}
    \item Any substituent with a lone pair of electrons on the atom bonded to the ring can provide resonance stabilization to a sigma complex.
    \begin{figure}[h]
        \centering
        \includegraphics[width=0.9\textwidth]{17.6_Activating_Group_op.png}
        \caption{Activating, Ortho, Para-Directors}
        \label{17.6_Activating_Group_op}
    \end{figure}

\end{enumerate}

\subsection{Deactivating, Meta-Directing Substituents 20260331}

\begin{enumerate}
    \item Nitrobenzene is about 100,000 times less reactive than benzene toward electrophilic aromatic substitution.
    \item An electron-donating substituent activates primarily the ortho and para positions, and an electron-withdrawing substituent (such as a nitro group) deactivates primarily the ortho and para positions.
    \item \textbf{Meta-directors}, often called meta-allowing substituents, deactivate the meta position less than the ortho and para positions, allowing meta substitution.
    \item Nitro group is a strong deactivating group by considering its resonance forms. %pic
    \item deactivating substituents are groups with a positive charge (or a partial positive charge) on the atom bonded to the aromatic ring. As we saw with the nitro group, this positively charged atom repels any positive charge on the adjacent carbon atom of the ring. Of the possible sigma complexes, only the one corresponding to meta substitution avoids putting a positive charge on this ring carbon. %...
    \begin{figure}[h]
        \centering
        \includegraphics[width=0.9\textwidth]{17.6_Deactivating_Group_m.png}
        \caption{Deactivating, Meta-Directors}
        \label{17.6_Deactivating_Group_m}
    \end{figure}
\end{enumerate}

\subsection{Halogen Substituents: Deactivating, but Ortho, Para-Directing 20260331}

\begin{enumerate}
    \item The halobenzenes are exceptions to the general rules. Halogens are deactivating groups, yet they are ortho, para-directors.
    \item The unusual combination of properties can be explained by considering that
    \begin{enumerate}
        \item the halogens are strongly electronegative, withdrawing electron density from a carbon atom through the sigma bond (inductive withdrawal).
        \item the halogens have nonbonding electrons that can donate electron density through pi bonding (resonance donation).
    \end{enumerate}
    % not so good, reference
    \item 
    \begin{figure}[h]
        \centering
        \includegraphics[width=0.9\textwidth]{17.6_Directing_Effects.png}
        \caption{Directing Effects of Substituents}
        \label{17.6_Directing_Effects}
    \end{figure}
\end{enumerate}

\subsection{Effects of Multiple Substituents on Electrophilic Aromatic Substitution 20260331}

\begin{enumerate}
    \item If the groups reinforce each other, the result is easy to predict.
\end{enumerate}

\begin{example} % 17-12
    Predict the mononitration products of the following compounds.
    \begin{enumerate}[label=(\alph*)]
        \item \textit{o}-nitrotoluene
        \item \textit{m}-chlorotoluene
        \item \textit{o}-bromobenzoic acid
        \item \textit{p}-methoxybenzoic acid
        \item \textit{m}-cresol (\textit{m}-methylphenol)
        \item \textit{o}-hydroxyacetophenone
    \end{enumerate}
\end{example}

\begin{enumerate}[resume]
    \item When the directing effects of two or more substituents conflict, it is more difficult to predict where an electrophile will react. In many cases, mixtures result.
    \item When there is a conflict between an activating group and a deactivating group, the activating group usually directs the substitution. In other words, activating groups are usually stronger directors than deactivating groups.
    \item Separation of substituents into three classes, from strongest to weakest.
    \begin{enumerate}
        \item Powerful ortho, para-directors that stabilize the sigma complexes through resonance.
        \item Moderate ortho, para-directors
        \item All meta-directors.
    \end{enumerate}
    \begin{figure}[h]
        \centering
        \includegraphics[width=0.9\textwidth]{17.6_Substituents_Strength.png}
        \caption{Substituents Strength}
        \label{17.6_Substituents_Strength}
    \end{figure}
    \item If two substituents direct an incoming electrophile toward different reaction sites, the substituent in the stronger class predominates. If both are in the same class, mixtures are likely.
\end{enumerate}

\begin{example}
    Predict the mononitration products of the following aromatic compounds.
    \begin{enumerate}[label=(\alph*)]
        \item \textit{p}-methylanisole
        \item \textit{m}-nitrochlorobenzene
        \item \textit{p}-chlorophenol
        \item \textit{m}-nitroanisole
    \end{enumerate}
\end{example}

\begin{example}
    \textit{Biphenyl} is two benzene rings joined by a single bond. The site of substitution for a biphenyl is determined by (1) which phenyl ring is more activated (or less deactivated), and (2) which position on that ring is most reactive, using the fact that a phenyl substituent is activating and ortho, para-directing.
    \begin{enumerate}[label=(\alph*)]
        \item Use resonance forms of a sigma complex to show why a phenyl substituent should be ortho, para-directing.
        \item Predict the mononitration products of the following compounds.
    \end{enumerate}
\end{example}

\subsection{The Friedel–Crafts Alkylation 20260331 20260402}

\begin{enumerate}
    \item In the presence of Lewis acid catalysts such as aluminum chloride (\ch{AlCl3}) or ferric chloride (\ch{FeCl3}), alkyl halides were found to alkylate benzene to give alkylbenzenes. This useful reaction is called the \textbf{Friedel–Crafts alkylation}.
    \begin{figure}[h] % title
        \centering
        \includegraphics[width=0.9\textwidth]{17.10_Friedel–Crafts_Alkylation.png}
        \caption{Friedel–Crafts alkylation}
        \label{17.10_Friedel–Crafts_Alkylation}
    \end{figure}
\end{enumerate}

\begin{mechanism}
    Friedel–Crafts Alkylation
\end{mechanism}

\begin{example} % 17-15
    Propose products (if any) and mechanisms for the following \ch{AlCl3}-catalyzed reactions:
    \begin{enumerate}[label=(\alph*)]
        \item chlorocyclohexane with benzene
        \item methyl chloride with anisole
        \item 3-chloro-2,2-dimethylbutane with isopropylbenzene
    \end{enumerate}
    
    \noindent \textbf{SOLUTION} \quad
    (c) rearrangement
\end{example}

\begin{enumerate}[resume]
    \item Friedel–Crafts Alkylation Using Other Carbocation Sources
    \begin{enumerate}
        \item Alkenes are protonated by \ch{HF} to give carbocations.
        \item Alcohols commonly form carbocations when treated with Lewis acids such as boron trifluoride (\ch{BF3}).
    \end{enumerate}
    \item Limitations of the Friedel–Crafts Alkylation
    \begin{enumerate}
        \item Friedel–Crafts reactions work only with \textbf{benzene}, \textbf{activated benzene derivatives}, and \textbf{halobenzenes}. They fail with strongly deactivated systems such as nitrobenzene, benzenesulfonic acid, and phenyl ketones.
        
        We can get around this limitation by adding the deactivating group or changing an activating group into a deactivating group \textit{after} the Friedel–Crafts step.

        \item The Friedel–Crafts alkylation is susceptible to carbocation rearrangements. Only certain alkylbenzenes can be made using the Friedel–Crafts alkylation.

        \item Because alkyl groups are activating substituents, the product of the Friedel–Crafts alkylation is \textit{more reactive} than the starting material. \textbf{Multiple alkylations} are hard to avoid.
        
        The problem of overalkylation can be minimized by using a large excess of benzene.
    \end{enumerate}
\end{enumerate}

\begin{example} % solved 17-2
    Devise a synthesis of \textit{p}-nitro-\textit{tert}-butylbenzene from benzene.
\end{example}

\begin{example} % 17-7
    Predict the products (if any) of the following reactions.
    \begin{enumerate}[label=(\alph*)]
        \item (excess) bezene + isobutyl chloride + \ch{AlCl3}
    \end{enumerate}
\end{example}

\begin{example}
    Which reactions will produce the desired product in good yield? You may assume that aluminum chloride is added as a catalyst in each case. For the reactions that will not give a good yield of the desired product, predict the major products.

    
    \begin{enumerate}
        \item 
    \end{enumerate}
\end{example}

\subsection{The Friedel–Crafts Acylation 20260402}

\begin{enumerate}
    \item An \textbf{acyl group} is a carbonyl group with an alkyl group attached. Acyl groups are named systematically by dropping the final -e from the alkane name and adding the \textit{-oyl} suffix.
    \item An \textbf{acyl chloride} is an acyl group bonded to a chlorine atom. Acyl chlorides are made by reaction of the corresponding carboxylic acids with thionyl chloride. Therefore, acyl chlorides are also called \textbf{acid chlorides}.
    \item In the presence of aluminum chloride, an acyl chloride reacts with benzene (or an activated benzene derivative) to give a phenyl ketone: an \textit{acylbenzene}.
    \item \textbf{Friedel–Crafts acylation} is analogous to the Friedel–Crafts alkylation, except that the reagent is an acyl chloride instead of an alkyl halide and the product is an acylbenzene (a “phenone”) instead of an alkylbenzene.
    \begin{figure}[h]
        \centering
        \includegraphics[width=0.9\textwidth]{17.10_Friedel–Crafts_Acylation.png}
        \caption{Friedel–Crafts acylation}
        \label{17.10_Friedel–Crafts_Acylation}
    \end{figure}
\end{enumerate}

\subsubsection{Mechanism of Acylation}

\begin{enumerate}
    \item The mechanism of Friedel–Crafts acylation (shown next) resembles that for alkylation, except that the electrophile is a resonance-stabilized acylium ion. The acylium ion reacts with benzene or an activated benzene derivative via an electrophilic aromatic substitution to form an acylbenzene.
\end{enumerate}

\begin{mechanism}
    Friedel–Crafts Acylation

    Friedel–Crafts acylation is an electrophilic aromatic substitution with an acylium ion acting as the electrophile.
\end{mechanism}

\begin{enumerate}[resume]
    \item Comparison of Friedel–Crafts Alkylation and Acylation
\end{enumerate}

\subsubsection{The Clemmensen Reduction: Synthesis of Alkylbenzenes}

\begin{enumerate}
    \item We use the Friedel–Crafts acylation to make the acylbenzene, then we reduce the acylbenzene to the alkylbenzene using the \textbf{Clemmensen reduction}: treatment with aqueous HCl and amalgamated zinc (zinc treated with mercury salts).
    \item The reagents and conditions for the Clemmensen reduction are similar to those used to reduce a nitro group to an amine.
\end{enumerate}

\subsubsection{The Gatterman–Koch Formylation: Synthesis of Benzaldehydes}

\begin{enumerate}
    \item We cannot add a formyl group to benzene by Friedel–Crafts acylation in the usual manner. Formyl chloride is unstable and cannot be bought or stored.
    \item Formylation can be accomplished by using a high-pressure mixture of carbon monoxide and HCl together with a catalyst consisting of a mixture of cuprous chloride (\ch{CuCl}) and aluminum chloride.
    \item The \textbf{Gatterman–Koch synthesis} is widely used in industry to synthesize aryl aldehydes.
\end{enumerate}

\subsection{Nucleophilic Aromatic Substitution 20260402}

\begin{enumerate}
    \item A nucleophile substitutes for a leaving group on an aromatic ring, and this class of reactions is called nucleophilic aromatic substitution.
    \item Electron-withdrawing substituents (such as nitro groups) \textit{activate} the ring toward nucleophilic aromatic substitution.
\end{enumerate}

\subsubsection{The Addition–Elimination Mechanism}

\begin{enumerate}
    \item Consider the reaction of 2,4-dinitrochlorobenzene with sodium hydroxide (shown next)...
\end{enumerate}

\begin{mechanism}
    Nucleophilic Aromatic Substitution (Addition–Elimination)
\end{mechanism}

\begin{enumerate}[resume]
    \item The resonance forms shown in the mechanism box illustrate how nitro groups ortho and para to the halogen help to stabilize the intermediate (and the transition state leading to it).
\end{enumerate}

\subsubsection{The Benzyne Mechanism: Elimination–Addition}

\begin{enumerate}
    \item Elimination–addition mechanism, called the \textbf{benzyne mechanism} because of the unusual intermediate.
    \item Sodium amide (or sodium hydroxide in the Dow process) reacts as a base, abstracting a proton.
\end{enumerate}

\begin{mechanism}
    Nucleophilic Aromatic Substitution (Benzyne Mechanism)
    The benzyne mechanism (elimination–addition) is likely when the ring has no strong electron-withdrawing groups. It usually requires a powerful base or high temperatures.
    \begin{itemize}
        \item Step 1: Deprotonation adjacent to the leaving group gives a carbanion.
        \item Step 2: The carbanion expels the leaving group to give a “benzyne” intermediate.
        \item Step 3: The nucleophile attacks at either end of the reactive benzyne triple bond.
        \item Step 4: Reprotonation gives the product.
    \end{itemize}
\end{mechanism}

\subsection{Aromatic Substitutions Using Organometallic Reagents 20260407}

\subsection{Addition Reactions of Benzene Derivatives 20260407}

\subsection{Side-Chain Reactions of Benzene Derivatives 20260407}

\subsubsection{Permanganate Oxidation}

\subsubsection{Side-Chain Halogenation}

\begin{enumerate}
    \item Alkylbenzenes undergo free-radical halogenation much more easily than alkanes because abstraction of a hydrogen atom at a benzylic position gives a resonancestabilized benzylic radical.
\end{enumerate}

\subsubsection{Nucleophilic Substitution at the Benzylic Position}

\begin{enumerate}
    \item Fisrt-Order Reactions: First-order nucleophilic substitution requires ionization of the halide to give a carbocation. In the case of a benzylic halide, the carbocation is resonance-stabilized.
    \item Second-Order Reactions: 
\end{enumerate}

\subsection{Reactions of Phenols 20260407}

\subsubsection{Oxidation of Phenols to Quinones}

\begin{enumerate}
    \item Phenols undergo oxidation, but they give different types of products from those seen with aliphatic alcohols. Chromic acid oxidation of a phenol gives a conjugated 1,4-diketone called a quinone.
\end{enumerate}

\subsubsection{Electrophilic Aromatic Substitution of Phenols}

\begin{enumerate}
    \item Phenols are highly reactive substrates for electrophilic aromatic substitution because the nonbonding electrons of the hydroxyl group stabilize the sigma complex formed by attack at the ortho or para position.
    \item the hydroxyl group is strongly activating and ortho, para-directing. Phenols are excellent substrates for halogenation, nitration, sulfonation, and some Friedel–Crafts reactions.
\end{enumerate}

\begin{example} % 17-43
    1,4-Benzoquinone is a good Diels–Alder dienophile. Predict the products of its reaction with
\end{example}

% HOMEWORK:17-44->47

\begin{example} % 17-47
    Predict the major products of bromination of the following compounds, using and FeBr3 in the dark.
    \begin{enumerate}[label=(\alph*)]
        \item 
    \end{enumerate}
\end{example}

\begin{example} % 17-51 ***
    Give the structures of compounds A through H in the following series of reactions.
\end{example}

\begin{example}
    A student added 3-phenylpropanoic acid to a molten salt consisting of a 1:1 mixture of \ch{NaCl} and \ch{AlCl3} maintained at 170 °C. After 5 minutes, he poured the molten mixture into water and extracted it into dichloromethane. Evaporation of the dichloromethane gave a 96\% yield of the product whose spectra follow. The mass spectrum of the product shows a molecular ion at $ m/z $ 132. What is the product?

    \noindent \textbf{SOLUTION} \quad indanone % 茚满酮
\end{example}

\begin{example}
    The following compound reacts with a hot, concentrated solution of NaOH (in a sealed tube) to give a mixture of two products. Propose structures for these products, and give a mechanism to account for their formation.

    \noindent \textbf{SOLUTION} \quad Considering benzyne
\end{example}

\begin{example} % 17-59
    \textit{Bisphenol} A is an important component of many polymers, including polycarbonates, polyurethanes, and epoxy resins. It is synthesized from phenol and acetone with HCl as a catalyst. Propose a mechanism for this reaction.
\end{example}

\begin{example} % 17-62
    When anthracene is added to the reaction of chlorobenzene with concentrated NaOH at 350 °C, an interesting Diels–Alder adduct of formula \ch{C20H14} results. The proton NMR spectrum of the product shows a singlet of area 2 around $ \delta \ 3 $ and a broad singlet of area 12 around $ \delta \ 7 $. Propose a structure for the product, and explain why one of the aromatic rings of anthracene reacted as a diene.

    \noindent \textbf{SOLUTION} \quad NaOH and \qty{350}{\degreeCelsius}, considering benzyne
\end{example}

\begin{example} % 17-64
    Phenol reacts with three equivalents of bromine in (in the dark) to give a product of formula When this product is added to bromine water, a yellow solid of molecular formula precipitates out of the solution. The IR spectrum of the yellow precipitate shows a strong absorption (much like that of a quinone) around Propose structures for the two products.
\end{example}

\begin{example} % 17-66
    A graduate student tried to make \textit{o}-fluorophenylmagnesium bromide by adding magnesium to an ether solution of \textit{o}-fluorobromobenzene. After obtaining puzzling results with this reaction, she repeated the reaction by using as solvent some tetrahydrofuran that contained a small amount of furan. From this reaction, she isolated a fair yield of the compound that follows. Propose a mechanism for its formation.
\end{example}

\begin{example} % 17-69
    The antioxidants BHA and BHT are commonly used as food preservatives. Show how BHA and BHT can be made from phenol and hydroquinone.
\end{example}